
\section{Introduction}

\textit{Data-Driven Computational Mechanics} (DDCM) constitutes the most recent paradigm shift in continuum mechanics, with certainly broad applicability. Specifically for the reaction-diffusion problem, DDCM seeks to find the concentration gradient and the associated flux vector field that best match an empirical extension of gradient-flux pairs. This is achieved by formulating an optimization problem where discrete and continuous fields are simultaneously sought to minimize their discrepancy, provided the continuous fields satisfy the physical conservation and geometrical compatibility laws as well as the Second Law of Thermodynamics.

A basic question that arises is the existence of optimal continuous fields when the discrete fields are held constant. This ``optimization subproblem,'' which is a continuous problem, can be analyzed and solved computationally using standard methods from the linear theory of \textit{Partial Differential Equations} (PDEs). Addressing this specific subproblem is the primary objective of this work. To set the stage for our approach, we first review key related works, differentiate our methodology from existing literature and state the novelty of our present contribution.

DDCM is a highly active research area whose origins date back less than ten years. Its core methodologies have predominantly been established within solid and structural mechanics, as exemplified by its application to the static analysis of multi-bar structures \cite{Kirchdoerfer2016}. Similar ideas led to the development of a data-driven magnetostatic solver \cite{Galetzka2020}. This initial approach was later modified to improve robustness in the presence of constitutive data outliers \cite{Kirchdoerfer2017} and subsequently extended to the dynamic analysis of structures featuring configuration-independent mass matrices \cite{Kirchdoerfer2018}. In the interest of full rigor, we note that the question of existence of solutions has also been addressed in the context of data-driven elasticity by several studies \cite{Conti2018,Conti2020,Gebhardt2026}.
Other research explores a hybrid strategy involving the offline reconstruction of a smooth constitutive manifold from empirical data. This concept was first implemented for time-independent analyses \cite{Gebhardt:2020a}, then adapted to solve time-dependent problems \cite{Gebhardt:2020b}. The approach has been further generalized to tackle challenging \textit{Mathematical Programs with Complementarity Constraints} (MPCCs), particularly relevant to contact mechanics simulations \cite{Gebhardt2024a}. Recently in  \cite{Gebhardt2024b}, the structural properties of the DDCM problem discretized by means of the \textit{Finite Element Method} (FEM) is investigated and its global solvability is established. Based on the same principles, DDCM is applied to discrete electrical circuits in \cite{Gebhardt2025}. Remarkably, the present work addresses a significant gap, as there is only a single antecedent of DDCM in the context of the diffusion-reaction problem, developed previously by some of the authors of this contribution \cite{Balbao2025} where a \textit{Galerkin/Least-Squares} (GLS) method for the stabilization of the resulting finite element discretization is proposed.

All previously introduced works that numerically solve the discrete-continuous optimization problem in DDCM rely on computing a solution to a ``optimization subproblem''. This is because those works adopt an \textit{Alternating Direction Method} (ADM). This work constitutes a crucial major improvement towards applying DDCM in the diffusion-reaction problem. Moreover, our analysis considers two or more spatial dimensions in contrast to contributions that focus on problems defined on closed intervals of the real line \cite{Gebhardt:2020a, Gebhardt:2020b, Gebhardt2024a, Gebhardt2024b}. We introduce novel finite element formulations and establish their well-posedness for the chosen discrete spaces. The formulations proposed are consistent, and therefore convergence is basically a consequence of stability, although we shall not pursue a convergence analysis of these formulations. The required stability properties are achieved by directly approximating the underlying saddle-point problem, an approach related to established methodologies such as those found in \cite{Masud2002, Bochev2006, Badia2010, Burman2023}.

Starting with the continuous setting, we present the primal and dual formulations for the PDE-constrained optimization problem and establish their well posedness. In addition and by leveraging the necessary regularity of the problem's functional setting, we successfully obtain the minimal irreducible forms of these two formulations. Then, we step into the discrete setting by adopting a finite element discretization. To ensure stability for both the discretized primal and dual problems when using continuous arbitrary approximations, we adopt the \textit{Variational MultiScale} (VMS) framework to systematically derive the necessary stabilization terms \cite{Hughes1995,Hughes1998,Codina2017}. In particular, we analyze some approximation's general properties and establish the consistency and well posedness of the resulting schemes when considering both the \textit{Algebraic Sub-Grid Scale} (ASGS) and \textit{Orthogonal Sub-Grid Scale} (OSGS) methods~\cite{codina2000, codina2002}. We show that one can move from the primal to the dual formulation just by changing appropriately the design of the stabilization terms, as in \cite{Badia2009} for Darcy's problem or in \cite{Badia2014} for the wave equation. We then stress-test the proposed discrete approximations and numerically investigate the crucial $\inf$-$\sup$ constant.  

The remainder of this article is structured as follows. Section 2 briefly establishes the context by reviewing the DDCM problem and derives through two main simplifications the PDE-constrained optimization problem targeted by the present contribution. Then, it culminates, in the continuous setting, with the formulation and analysis of the associated primal and dual forms. Section 3 formulates stable finite element approximations for the primal and dual forms and establishes their well posedness. Section 4 details the numerical implementation, realized using the Firedrake finite element platform, and presents a series of progressively sophisticated examples to illustrate the proposed setting's properties. Section 5 provides concluding remarks, discusses limitations and suggests possible avenues for future work.
